\section{Related Work}
\label{sec:related}

\paragraph{Fixed-budget KV-cache compression.} A large body of work bounds the
key--value cache to a fixed budget by discarding or merging entries along the
\emph{token} axis. StreamingLLM~\citep{streamingllm} keeps only a recent window
plus a few initial attention-sink tokens; H2O~\citep{h2o} evicts all but the
recent and heavy-hitter tokens ranked by accumulated attention mass;
SnapKV~\citep{snapkv} predicts, from a short observation window, which positions
each head will attend to and prunes the rest before generation. Others vary the
budget \emph{across} layers---PyramidKV~\citep{pyramidkv} allocates more cache to
lower layers following an information-funnel observation, and
LCKV~\citep{lckv} condenses the cache to a few layers whose KV the remaining
layers share---while Activation Beacon~\citep{actbeacon} learns, through
continued training, to compress per-layer activations into a small set of beacon
tokens. Activation Beacon is one instance of a broader
\emph{compress-then-inject} family---Gist tokens~\citep{gist},
AutoCompressor~\citep{autocompressor}, ICAE~\citep{icae}, and
xRAG~\citep{xrag}---that encodes context into a few soft tokens or vectors and
injects them into the model; CoMem does the opposite, keeping the raw hidden and
partitioning on \emph{depth} rather than compressing the token count.
CompressKV~\citep{compresskv} adds a semantic retrieval-head criterion to
token eviction and is orthogonal to us, as it still operates on the token axis.
The common denominator is that all of these methods retain (a compressed form of)
the \emph{full-depth} KV of the kept tokens; CoMem instead caches a single
mid-depth hidden state and reconstructs the upper layers on demand. We compare
head-to-head against StreamingLLM under an identical read budget in
\S\ref{sec:experiments}.

\paragraph{Retrieval-augmented long context.} A second line augments a bounded
working set with a retrieval step that surfaces the relevant distant context.
Landmark Attention~\citep{landmark} trains a landmark token per block and lets
attention select blocks through it, while FoT/LongLLaMA~\citep{longllama} attach
an external $(k,v)$ memory and use contrastive training to keep retrieved keys
distinguishable; Memorizing Transformers~\citep{memtransformers} similarly cache
$(k,v)$ pairs at a \emph{single} layer and augment that layer's attention with an
approximate $k$NN lookup. On the serving side, CacheBlend~\citep{cacheblend} and
RAGCache~\citep{ragcache} precompute and reuse the KV of retrieved chunks,
selectively recomputing a few tokens to restore cross-attention quality;
Block-Attention~\citep{blockattn} splits retrieved documents into blocks that
each compute their KV independently, re-encodes positions, and fine-tunes the
model to tolerate the block-wise reused cache; and
InfLLM~\citep{infllm} stores distant context as memory units and looks up
token-relevant units inside attention, training-free. In every case the
retrieved unit is the \emph{full-layer} KV of a block, and retrieval is either
folded into attention (differentiable, per-step) or organized as a KV store to
be attended. Closest in spirit to retrieving over stored \emph{hidden} rather
than KV, Unlimiformer~\citep{unlimiformer} offloads cross-attention to a single
$k$NN index over an encoder--decoder's top-layer hidden states, and
CEPE~\citep{cepe} encodes chunks in parallel with a small auxiliary encoder that
a frozen decoder cross-attends to. Both consume cached hidden states, but they
feed them into (cross-)attention and never resume a forward pass, and CEPE adds a
separate encoder to train. CoMem departs from all of these: it retrieves with an
external, non-differentiable, one-shot BM25 pass over chunk text, and---to our
knowledge uniquely---it \emph{recomputes} the upper layers of its own backbone
from a cached mid-layer hidden rather than attending to cached KV or hidden.

\paragraph{Fixed-size / architectural long-context memory.} A third line encodes
long-range information into a memory whose size does not grow with context.
MemoryLLM~\citep{memoryllm} maintains a fixed-size latent memory pool inside the
transformer that is self-updated from new text and survives many updates; we
also compare against it as a same-class rival in \S\ref{sec:experiments}.
RecursiveSummarizing~\citep{recsum} keeps a text-space memory by recursively
rewriting a natural-language summary, and YOCO~\citep{yoco} redesigns the model
as a decoder--decoder that caches a single global KV once and shares it across
upper cross-decoder layers, reaching near-perfect million-token retrieval (we
cite its numbers as an anchor rather than reproducing it). What unifies these
methods is that the memory is either baked into the parameters or requires
retraining the architecture from scratch. CoMem keeps the backbone stock, adds
only a light self-distilled LoRA, and stores an external, non-evolving hidden
cache that is read by recomputation.

\paragraph{Caching intermediate activations and recomputing upper layers (closest).}
Nearest to us are methods that cache a partial forward state and reconstruct the
rest. HCache~\citep{hcache} restores an evicted KV cache in a serving system by
storing a mid-layer activation and recomputing the layers above it, but it is
post-hoc, untrained, and does no retrieval. KV-Direct~\citep{kvdirect}
(``The Residual Stream Is All You Need'') shows the KV cache is a deterministic
projection of the residual stream and recomputes the \emph{full-depth} KV from a
cached residual, training-free and numerically exact, but it keeps every token so
its working set still grows with context. KV-CAT~\citep{kvcat} induces cache
compressibility through continued pretraining on the token/slot axis---we cite it
as prior art and explicitly do \emph{not} claim its train-for-compressibility
mechanism as ours, since CoMem partitions on the \emph{depth} axis and works
zero-training. LLoCO~\citep{lloco} is the closest \emph{pipeline}-level match: it
processes long context offline, compresses it, adapts a LoRA so the model can
read the compressed representation, and retrieves the relevant compressed context
at query time. We do \emph{not} claim the offline$+$LoRA$+$retrieval framework as
original---LLoCO owns that framing; our departure is \emph{what} is stored and
\emph{how} it is read. LLoCO compresses context into lossy summary embeddings
prepended as a soft prompt and never re-runs the backbone over the original
tokens, whereas CoMem stores the \emph{raw} layer-$j$ hidden and \emph{resumes}
the forward from layer $j$---a depth knob rather than a compression ratio---with
a self-distillation target specific to that depth partition.
StagesOfInference~\citep{stagesofinference} provides the
depth-semantics rationale---distinct functional stages emerge along the layer
stack---that motivates treating depth as the partition axis. Early-exit methods
such as LayerSkip~\citep{layerskip} and CALM~\citep{calm} exploit the same
``layers have distinct roles'' premise from the opposite direction: they skip the
\emph{upper} layers for easy tokens, whereas CoMem caches the output of the
\emph{lower} layers and always recomputes the upper ones with the query present.
CoMem occupies a
unique intersection of these ideas: layer-\emph{partial} recompute (only
layers $[j{:}L]$) driven by external BM25 retrieval, a read length that is
constant in context length, a self-distillation step that pushes the cacheable
split depth deeper, and the split $j$ exposed as a single RAG$\leftrightarrow$closed-book
knob.

\paragraph{Novelty, positioned honestly.} CoMem is, to our knowledge, the first
long-context method to partition the transformer on the \emph{depth} axis---%
caching a single mid-layer hidden per chunk and recomputing only the upper layers
$[j{:}L]$---driven by external retrieval into a constant-length resumed forward
and stabilized by self-distillation; its nearest neighbors either recompute
without retrieval (HCache, KV-Direct, Block-Attention) or retrieve without a
layer-partial resume (LLoCO, InfLLM, Unlimiformer). We are deliberate about which
primitives are new. The depth-axis partition with layer-partial recompute is the
genuinely new ingredient: the only prior work that caches a mid-layer activation
and recomputes above it is HCache, a training-free serving-cache restore with no
retrieval and no notion of depth as a knob. Feeding retrieval into a
constant-length \emph{resumed} forward is a new combination of otherwise
well-precedented parts. The top-level offline$+$LoRA$+$retrieval recipe we
explicitly do \emph{not} claim as novel---LLoCO owns that framing, and the
finetune-to-tolerate-reuse idea is shared with Block-Attention; our
differentiator inside it is the self-distillation target (a full-context teacher
distilled into the resume-from-$j$ student on the teacher's top-support tokens),
which is specific to the depth partition and, as far as our survey found, not
used elsewhere.
